\subsection*{Determinant calls rescue prediction on unseen lineages}

The two-group pattern above is a statement about features as much as about lineages. For the three
drugs that lose most under the holdout, resistance is usually conferred by an acquired gene, and a
reference pan-genome catalogue represents such a gene only if a strain carrying it was in the
reference. The pre-registered test of that reading replaces the catalogue with each strain's own
determinant calls.

With ST-grouped five-fold cross-validation, the catalogue alone reached a mean \nap{} of 0.450 over
the three drugs with acquired determinants (CTX, CAZ, GEN). Adding the determinant calls raised it
to 0.726, a difference of $+0.276$ [0.178, 0.328], while ciprofloxacin, whose determinants are
point mutations the catalogue already carries, did not move ($-0.004$; Figure~\ref{fig:rescue}).
The rescue held with intermediates excluded ($+0.293$), across all database versions ($+0.286$), on
the country splits (UK $\rightarrow$ Norway 0.249 $\rightarrow$ 0.429; Norway $\rightarrow$ UK
0.543 $\rightarrow$ 0.793), on the temporal split (0.449 $\rightarrow$ 0.670), and inside each
country separately (United Kingdom $+0.256$ [0.119, 0.457]; Norway $+0.308$ [0.190, 0.363]). The
calls alone scored as well as calls plus catalogue (0.730 against 0.726): for these drugs, on
unseen lineages, the catalogue added nothing.

The rescue is demonstrated on the 5,067 strains with both a catalogue profile and published calls,
and one pre-registered check of that restriction failed. Asked whether the catalogue-only collapse
is equally present in the 1,890 strains without calls, we found the catalogue scoring 0.641 there
against 0.449 in the called strains, outside the $\pm0.10$ tolerance set in advance. Most of the
uncalled strains (1,439) come from a single study, and within that stratum metadata alone predicts
the cephalosporin phenotypes well (CTX \nap{} 0.735), which points to study-specific structure; the
analysis was not designed to separate that from a genuine difference in the strains and we claim
neither. The consequence for the paper is that the rescue is established on the called subset.

\begin{figure}[t]\centering
\figbox{Figure 4. The determinant rescue}{6.5cm}{Panel A: \nap{} per drug on ST-grouped 5-fold for
catalogue, catalogue $+$ calls, calls only, and metadata. Panel B: catalogue against catalogue $+$
calls (mean of CTX, CAZ, GEN) on every split, with the CIP control. Panel C: the same inside the
United Kingdom and inside Norway.}{results/w6\_rescue/paper\_ri/metrics.csv, deltas.csv;
results/w10\_coverage\_rule/part\_c.csv}
\caption{\textbf{Isolate-specific determinant calls rescue prediction on unseen lineages.}
\todo{caption}}\label{fig:rescue}
\end{figure}

\subsection*{Assembly size is a within-lineage proxy for accessory content}

Assembly statistics held more signal than metadata on every split that moved populations, which
raises the question of what they are measuring. Inside sequence types, assembly length correlated
with the number of acquired resistance genes (within-ST Spearman $\rho = 0.325$; 4,033 strains, 62
STs) and ranked resistant above susceptible strains of the same type (within-ST concordance
0.552--0.676 across the seven drugs, where an ST-constant score would give exactly 0.5).

A batch artefact would explain this as well as biology, so we conditioned on the submitting study
too. Within (ST $\times$ study) cells the concordance of assembly length over the
acquired-determinant drugs fell only from 0.613 to 0.592, and its correlation with acquired-gene
count from 0.325 to 0.321 (Figure~\ref{fig:assembly}). The signal is carried by content statistics
rather than technical ones: contig N50 stayed below 0.5 for every drug (0.383--0.490) and
correlated negatively with acquired-gene count ($\rho = -0.185$), as expected if larger, more
fragmented assemblies carry more plasmid DNA. The same pattern appeared in \textit{Campylobacter}
for tetracycline, whose tet(O) gene is commonly plasmid-borne \cite{iovine2013campylobacter}
(within-cluster concordance 0.832), and not for the chromosomal point-mutation drugs
(0.450--0.549). The mean over its drugs, 0.595, met the pre-registered threshold while being
carried almost entirely by one drug.

This is real biology inside a collection and it is worthless outside one. On the external cohort a
model on assembly statistics alone reached a mean \nap{} of 0.045.

\begin{figure}[t]\centering
\figbox{Figure 5. The assembly proxy}{7cm}{Panel A: within-ST and within-(ST$\times$study)
concordance per assembly statistic and drug, as a heatmap, content statistics against technical
ones. Panel B: assembly length against acquired-gene count, centred within ST. Panel C:
\textit{Campylobacter} within-cluster concordance of length per drug, tetracycline against the
point-mutation drugs. Panel D: assembly-arm \nap{}, internal against
external.}{results/w12\_shortcut\_mechanism/*.csv;
results/w13\_campylobacter/assembly\_within\_cluster.csv;
results/w8\_external/primary\_ri/metrics.csv}
\caption{\textbf{Assembly statistics track accessory gene content within lineages and do not
transfer.} \todo{caption}}\label{fig:assembly}
\end{figure}

\subsection*{Frozen models transfer to an external cohort}

The external cohort differs from the training collection in every respect we could measure: 68\% of
its isolates are from the United States, 57\% from food or animal sources, and its median
collection year is 2019, against a two-country research collection spanning 2001--2017
(Table~\ref{tab:cohorts}). Gradient boosting on all determinant calls, trained on our strains and
applied unchanged, reached a mean \nap{} of 0.807 over six drugs, against 0.045 for the assembly
arm: a pre-registered difference of $+0.762$ [0.676, 0.804] over 62 BioProjects
(Figure~\ref{fig:external}). It held across all database versions ($+0.773$) and with intermediates
excluded ($+0.753$). Dropping each of the nine BioProjects with at least 100 labelled isolates in
turn left the difference between 0.723 and 0.777, with no interval below 0.643 and no overlap with
zero. No external isolate matched any identifier of a training strain.

\begin{figure}[t]\centering
\figbox{Figure 6. External validation}{6.5cm}{Panel A: \nap{} per drug for all calls (XGBoost),
acquired calls only, L1 logistic regression, the determinant rule and the assembly arm, on 10,232
external isolates. Panel B: leave-one-BioProject-out estimates of the difference over the assembly
arm and over the rule.}{results/w8\_external/primary\_ri/metrics.csv;
results/w14\_clinical\_linear/b\_metrics.csv; results/w11\_external\_robustness/leave\_one\_out.csv}
\caption{\textbf{Frozen models on an external cohort.} \todo{caption}}\label{fig:external}
\end{figure}

\subsection*{A sparse linear model on determinant calls is enough where a model helps}

Once the determinants are in the features, the baseline a model has to beat is not another model but
the obvious rule: call a strain resistant when it carries a known determinant. On unseen lineages a
learned model does beat it, by $+0.122$ [0.075, 0.197] mean \nap{} over CTX, CAZ and GEN, with every
per-drug interval above zero.

The learning that matters is not the flexible part. An $L_1$ logistic regression on the same
features did as well as gradient boosting internally (difference $-0.011$ [$-0.026$, 0.007], inside
the pre-registered equivalence margin of $\pm0.05$) and beat the rule by $+0.137$ [0.097, 0.203]
(Figure~\ref{fig:rule}). Externally the ordering reverses. Gradient boosting beat the rule by only
$+0.065$ [0.005, 0.140] and lost to it for GEN and CTX; the linear model beat the rule by $+0.104$
[0.060, 0.166], was at or above it for every drug, and was ahead of gradient boosting on six of
seven drugs (mean $+0.038$ [0.019, 0.063]; by the pre-registered margin of 0.05 this comparison is
inconclusive, and we state the direction rather than claim the effect). The two drugs on which the
rule beat gradient boosting are therefore failures of that model rather than of learning from
determinant calls. External mean \nap{}: linear 0.849, gradient boosting 0.785, rule 0.755.

Label heterogeneity hides part of the margin. Re-labelling the external isolates with EUCAST
breakpoints from their MICs raised the margin of gradient boosting over the rule from $+0.065$ to
$+0.123$ [0.081, 0.174]. Because EUCAST labels exist for more isolates than submitter calls, we
repeated the comparison on the 35,367 isolate--drug pairs carrying both: the margin rose from
$+0.064$ [$-0.009$, 0.159] to $+0.101$ [0.066, 0.158]. That comparison was exploratory, not
pre-registered.

The advantage of a model is not general, and the second species is where it ends. In
\textit{Campylobacter}, where resistance to the four eligible drugs is explained by gyrA T86I,
tet(O) and 23S rRNA mutations \cite{zhao2016campylobacter}, the rule reached a mean \nap{} of 0.928
and gradient boosting 0.922; the pre-registered difference was $-0.006$ [$-0.017$, 0.004], a
failure of the hypothesis that a model helps. Where the determinants are nearly sufficient, there
is nothing left for a model to learn.

\begin{figure}[t]\centering
\figbox{Figure 7. Model versus rule}{7cm}{Forest plot of (model $-$ rule) mean \nap{} with
intervals: \Ecoli{} unseen lineages (XGBoost, L1), \Ecoli{} external with submitter labels
(XGBoost, L1), external with EUCAST labels on identical isolates, and \textit{Campylobacter}.
Inset: per-drug external differences for XGBoost and L1.}{results/w10\_coverage\_rule/deltas.csv;
results/w14\_clinical\_linear/b\_deltas.csv; results/w13\_mic/m1\_*.csv;
results/w13\_campylobacter/deltas.csv}
\caption{\textbf{When a learned model beats a determinant rule.} \todo{caption}}\label{fig:rule}
\end{figure}

\subsection*{Very major errors remain the unsolved problem}

At the fixed threshold, with EUCAST labels, categorical agreement over the six external drugs was
0.914 for gradient boosting, 0.938 for the linear model and 0.942 for the rule, and major error
rates were 2.7\%, 1.7\% and 2.7\%. Very major error rates were 25.1\%, 18.7\% and 14.4\%: no
approach met the regulatory benchmark, and only the rule for ciprofloxacin did so for a single drug
(2.6\%, upper limit 3.5\%; Table~\ref{tab:clinical}, Figure~\ref{fig:clinical}). Paired and
clustered by BioProject, the very major error rate of gradient boosting exceeded the rule's by 10.7
points [3.2, 18.1] and the linear model's by 4.3 [0.6, 8.4], while major error rates were
indistinguishable. Amoxicillin--clavulanic acid was the worst drug for every approach (very major
error 35--50\%), as the uncatalogued variation in its resistance genes predicts
\cite{lipworth2024uncatalogued}.

So the learned models rank resistant isolates better than the rule and, at the pre-stated
threshold, call fewer of them. This is the same operating-point problem the internal threshold sweep
found, and we did not move the threshold after seeing it.

Errors were more frequent near the breakpoint than far from it, but by less than the pre-registered
factor of two (ratio 1.64 for gradient boosting, 1.87 for the rule), so that hypothesis failed as
stated. The per-drug picture is more informative than the pooled ratio: for ciprofloxacin the
model's errors were almost all within one dilution of the breakpoint, while for gentamicin 18.1\%
were far from it, against 2.7\% for the rule.

\begin{figure}[t]\centering
\figbox{Figure 8. Clinical error rates}{6cm}{Per drug: very major and major error rates with
Clopper--Pearson intervals for XGBoost, L1 and the rule, with the regulatory benchmarks as
reference lines. Second panel: error rate near against far from the EUCAST
breakpoint.}{results/w14\_clinical\_linear/a\_clinical.csv; results/w13\_mic/m2\_per\_drug.csv}
\caption{\textbf{Categorical errors on the external cohort.} \todo{caption}}\label{fig:clinical}
\end{figure}
